\section{Tổng quan}
%-------------------------------------------------
% Frame 1
%-------------------------------------------------
\begin{frame}{Tại sao cần kiểm định giả thuyết?}

\begin{block}{Quan điểm của Karl Popper}
Một giả thuyết khoa học phải có khả năng được kiểm chứng bằng dữ liệu thực nghiệm.
\end{block}

\begin{itemize}
\item Giả thuyết phải có khả năng bị bác bỏ.
\item Nếu không thể kiểm tra bằng dữ liệu thì đó không phải là giả thuyết khoa học.
\item Dữ liệu và thống kê đóng vai trò trung tâm trong phương pháp khoa học hiện đại.
\end{itemize}

\end{frame}

%-------------------------------------------------
% Frame 2
%-------------------------------------------------
\begin{frame}{Kiểm định giả thuyết trong thống kê cổ điển}

\begin{itemize}

\item Bắt đầu với giả thuyết không: $H_0$
\item Giả sử $H_0$ là đúng.

\item Thu thập dữ liệu từ thí nghiệm.

\item Nếu dữ liệu quan sát được quá bất thường dưới $H_0$ thì bác bỏ $H_0$.

\item Mức độ bằng chứng thường được đo bằng \textbf{p-value}.

\end{itemize}

\begin{block}{Ý tưởng chính}
Kiểm định giả thuyết có thể được xem như một dạng chứng minh phản chứng ngẫu nhiên.
\end{block}

\end{frame}

%-------------------------------------------------
% Frame 3
%-------------------------------------------------
\begin{frame}{Hạn chế của p-value}

\begin{itemize}

\item p-value truyền thống thường giả định kích thước mẫu cố định.

\item Nhiều ứng dụng hiện đại tạo dữ liệu theo cách tuần tự:
\begin{itemize}
    \item A/B Testing
    \item Thử nghiệm lâm sàng
    \item Hệ thống học máy
    \item Giám sát thời gian thực
\end{itemize}

\item Người nghiên cứu có thể liên tục kiểm tra kết quả trong quá trình thu thập dữ liệu.

\end{itemize}

\begin{alertblock}{Vấn đề}
Việc dừng thí nghiệm tùy ý (optional stopping) có thể làm mất tính hợp lệ của suy luận thống kê cổ điển.
\end{alertblock}

\end{frame}

%-------------------------------------------------
% Frame 4
%-------------------------------------------------
\begin{frame}{SAVI: Sequential Anytime-Valid Inference}

\begin{block}{Mục tiêu}
Xây dựng các phương pháp suy luận vẫn đúng tại mọi thời điểm dừng quan sát.
\end{block}

\begin{itemize}
\item Dữ liệu có thể được theo dõi liên tục.
\item Thí nghiệm có thể dừng ở bất kỳ thời điểm nào.
\item Vẫn kiểm soát được sai lầm loại I.
\item Không cần xác định trước kích thước mẫu cuối cùng.
\end{itemize}

\end{frame}

%-------------------------------------------------
% Frame 5
%-------------------------------------------------
\begin{frame}{Ý tưởng cá cược chống lại giả thuyết không}

\begin{itemize}

    \item SAVI diễn giải kiểm định giả thuyết như một trò chơi cá cược chống lại $H_0$.

    \item Một người chơi liên tục đặt cược dựa trên dữ liệu thu được.

    \item Nếu tài sản tăng nhanh, điều đó cho thấy dữ liệu không phù hợp với $H_0$.

\end{itemize}

\medskip

\[
\text{Tài sản tăng}
\Longrightarrow
\text{Bằng chứng chống lại } H_0
\]

\medskip

\begin{block}{Trực giác}
Bác bỏ $H_0$ tương đương với việc kiếm được tiền khi cá cược chống lại $H_0$.
\end{block}

\end{frame}

%-------------------------------------------------
% Frame 6
%-------------------------------------------------
\begin{frame}{Từ cá cược đến Game-Theoretic Statistics}

\begin{itemize}

\item Cách tiếp cận cá cược dẫn đến:
\begin{itemize}
    \item Test Martingale
    \item e-value
    \item e-process
\end{itemize}

\item Đồng thời tạo cầu nối với:
\begin{itemize}
    \item Lý thuyết thông tin
    \item Tiêu chuẩn Kelly
    \item Xác suất và tài chính
    \item Học trực tuyến và tối ưu hóa
\end{itemize}

\item Đây chính là nền tảng của \textbf{Game-Theoretic Statistics}.

\end{itemize}

\end{frame}

%====================================================
% SAVI - Frame 1
%====================================================
\subsection{Sequential Anytime-Valid Inference (SAVI)}

\begin{frame}{P-value trong thống kê cổ điển}

\begin{itemize}

\item Giả sử $P(n)$ là p-value được tính từ $n$ quan sát đầu tiên.

\item Với mức ý nghĩa $\alpha$, dưới giả thuyết không $H_0$:

\[
\Pr(P(n)\le \alpha)\le \alpha
\]

\item Điều này có nghĩa rằng xác suất sai lầm loại I được kiểm soát ở mức $\alpha$.

\item Tuy nhiên, kết quả trên chỉ đúng khi kích thước mẫu $n$ được xác định trước.

\end{itemize}

\begin{block}{Thông điệp chính}
P-value cổ điển chỉ đảm bảo tính hợp lệ khi số lượng dữ liệu được cố định trước khi thí nghiệm bắt đầu.
\end{block}

\end{frame}

%====================================================
% SAVI - Frame 2
%====================================================
\begin{frame}{Vấn đề của dữ liệu tuần tự}

\begin{itemize}

\item Trong thực tế, dữ liệu thường xuất hiện theo thời gian.

\item Nhà nghiên cứu có thể:
\begin{itemize}
\item Quan sát dữ liệu liên tục;
\item Kiểm tra p-value nhiều lần;
\item Quyết định dừng hoặc tiếp tục thí nghiệm.
\end{itemize}

\item Gọi $\tau$ là thời điểm dừng của thí nghiệm.

\item Khi đó ta quan tâm đến:

\[
\Pr(P(\tau)\le \alpha)
\]

\end{itemize}

\begin{alertblock}{Vấn đề}

Thông thường:

\[
\Pr(P(\tau)\le \alpha)\nleq \alpha
\]

\end{alertblock}

\end{frame}

%====================================================
% SAVI - Frame 3
%====================================================
\begin{frame}{Optional Stopping}

\begin{block}{Chiến lược phổ biến}
Tiếp tục thu thập dữ liệu cho đến khi p-value trở nên có ý nghĩa thống kê.
\end{block}

\[
\tau=
\min
\left\{
n\in\mathbb{N}:
P(n)\le \alpha
\right\}
\]

\begin{itemize}

\item Liên tục kiểm tra p-value.

\item Dừng ngay khi $P(n)\le \alpha$.

\item Khi đó xác suất bác bỏ sai $H_0$ có thể lớn hơn rất nhiều so với $\alpha$.

\end{itemize}

\begin{alertblock}{Kết quả}
Nếu tiếp tục đủ lâu, xác suất dương tính giả có thể tiến gần tới $1$.
\end{alertblock}

\end{frame}

%====================================================
% SAVI - Frame 4
%====================================================
\begin{frame}{Khoảng tin cậy cũng gặp vấn đề}

Giả sử

\[
(L^{(n)},U^{(n)})
\]

là khoảng tin cậy mức $(1-\alpha)$.

\vspace{0.3cm}

Khi $n$ cố định:

\[
\Pr
\Big(
\theta\in(L^{(n)},U^{(n)})
\Big)
\ge 1-\alpha
\]

\vspace{0.3cm}

Nếu tiếp tục thu thập dữ liệu cho đến khi

\[
L^{(n)}>0
\]

thì thời điểm dừng phụ thuộc vào dữ liệu.

\begin{alertblock}{Hệ quả}
Tính chất bao phủ của khoảng tin cậy không còn được đảm bảo.
\end{alertblock}

\end{frame}

%====================================================
% SAVI - Frame 5
%====================================================
\begin{frame}{Động lực ra đời của SAVI}

\begin{itemize}

\item P-value mất tính hợp lệ dưới optional stopping.

\item Khoảng tin cậy cổ điển cũng mất tính bao phủ.

\item Dữ liệu hiện đại thường được theo dõi liên tục:

\begin{itemize}
\item A/B testing;
\item Clinical trials;
\item Online learning;
\item Real-time monitoring.
\end{itemize}

\end{itemize}

\begin{block}{Câu hỏi}
Liệu có thể xây dựng một phương pháp suy luận vẫn đúng dù dừng ở bất kỳ thời điểm nào?
\end{block}

\end{frame}

%====================================================
% SAVI - Frame 6
%====================================================
\begin{frame}{Sequential Anytime-Valid Inference (SAVI)}

\begin{block}{Ý tưởng cốt lõi}
Một phương pháp được gọi là SAVI nếu nó vẫn duy trì tính hợp lệ tại mọi thời điểm dừng.
\end{block}

\begin{itemize}

\item Theo dõi dữ liệu liên tục.

\item Dừng hoặc tiếp tục thí nghiệm vì bất kỳ lý do nào.

\item Thực hiện peeking nhiều lần.

\item Vẫn kiểm soát sai lầm loại I.

\end{itemize}

\[
\text{Validity}
\Longrightarrow
\text{At Every Stopping Time}
\]

\end{frame}

%====================================================
% SAVI - Frame 7
%====================================================
\begin{frame}{Nền tảng toán học của SAVI}

Các công cụ cốt lõi của SAVI bao gồm:

\begin{itemize}

\item Test Martingales

\item e-values

\item e-processes

\end{itemize}

\vspace{0.3cm}

Các công cụ này cho phép:

\begin{itemize}

\item Đo lường bằng chứng thống kê theo thời gian

\item Duy trì tính hợp lệ dưới optional stopping

\item Hỗ trợ suy luận tuần tự và anytime-valid

\end{itemize}

\begin{block}{Phần tiếp theo}
Chúng ta sẽ bắt đầu từ ý tưởng cá cược và test martingale.
\end{block}

\end{frame}

%====================================================
% SAVI - Frame 8: The Lady Tasting Tea
%====================================================
\begin{frame}{Thí nghiệm: The Lady Tasting Tea (Fisher, 1920s)}

\begin{block}{Câu chuyện lịch sử}
Muriel Bristol tuyên bố bà có thể nhận biết được trà được rót vào sữa trước hay sữa được rót vào trà trước.
\end{block}

\begin{itemize}
    \item \textbf{Giả thuyết không $H_0$:} Muriel chỉ đang đoán mò (xác suất 50/50).
    \item \textbf{Thiết kế:} 8 cốc trà (4 loại này, 4 loại kia), yêu cầu phân loại.
\end{itemize}

\vspace{0.2cm}
\textbf{Kết quả dưới $H_0$:}
\begin{itemize}
    \item Nếu đúng hoàn toàn: p-value $= \frac{1}{70} \approx 0.014$ (Bác bỏ $H_0$).
    \item Nếu chỉ sai 1 ly: p-value $= \frac{17}{70} \approx 0.24$ (Không thể bác bỏ $H_0$).
\end{itemize}

\end{frame}

%====================================================
% SAVI - Frame 9: Hạn chế
%====================================================
\begin{frame}{Hạn chế của phương pháp cổ điển}

Từ thí nghiệm của Fisher, ta thấy các nhược điểm cốt lõi:

\begin{itemize}
    \item \textbf{Mất mát bằng chứng:} Chỉ một sai sót nhỏ cũng làm bằng chứng thống kê (p-value) suy giảm nghiêm trọng.
    \item \textbf{Tổng hợp một lần:} Toàn bộ bằng chứng chỉ được tính toán một lần duy nhất vào cuối thí nghiệm.
    \item \textbf{Kích thước mẫu cố định:} Không được dừng sớm nếu bằng chứng đã rõ ràng, cũng không được thu thập thêm dữ liệu nếu bằng chứng chưa đủ mạnh (vấn đề \textit{Optional Stopping}).
\end{itemize}

\end{frame}

%====================================================
% SAVI - Frame 10: The Lady Keeps Tasting Coffee
%====================================================
\begin{frame}{The Lady Keeps Tasting Coffee (Ramdas)}

Để khắc phục, Aaditya Ramdas đề xuất một phiên bản hiện đại:

\begin{block}{Thiết lập thí nghiệm mới}
Leila Wehbe cần phân biệt: Espresso rót vào sữa hay sữa rót vào Espresso. Cô có thể tiếp tục quan sát và dự đoán trong thời gian \textbf{tùy ý}.
\end{block}

\begin{itemize}
    \item \textbf{$H_0$:} Leila không có khả năng phân biệt (xác suất đoán đúng ở mỗi vòng chỉ là $1/2$).
    \item \textbf{Cách tiếp cận SAVI:} Không đếm số câu trả lời đúng để tính p-value, mà thiết kế một \textbf{trò chơi cá cược} giữa nhà thống kê và giả thuyết không.
\end{itemize}

\end{frame}

%====================================================
% SAVI - Frame 11: Ý tưởng cá cược
%====================================================
\begin{frame}{Kiểm định thông qua cá cược (Testing by Betting)}

\begin{block}{Cơ chế cá cược}
\begin{itemize}
    \item Ở mỗi dự đoán, Leila đặt cược một phần tài sản đang có.
    \item Đoán đúng $\rightarrow$ Tài sản tăng. Đoán sai $\rightarrow$ Tài sản giảm.
\end{itemize}
\end{block}

\textbf{Logic của phương pháp:}
\begin{itemize}
    \item Nếu $H_0$ ĐÚNG: Trò chơi là công bằng. Không có chiến lược nào giúp làm giàu một cách hệ thống.
    \item Nếu $H_0$ SAI: Tồn tại sự khác biệt, một chiến lược tốt sẽ tạo ra lợi nhuận.
\end{itemize}

\begin{alertblock}{Kết luận cốt lõi}
Tài sản tích lũy được trong trò chơi chính là \textbf{thước đo trực tiếp} của bằng chứng chống lại giả thuyết không.
\end{alertblock}

\end{frame}

%====================================================
% SAVI - Frame 12: Thiết lập toán học
%====================================================
\begin{frame}{Thiết lập toán học của quá trình cá cược}

\begin{itemize}
    \item \textbf{Tài sản ban đầu:} $L_0 = 1$.
    \item \textbf{Tỷ lệ đặt cược:} $\lambda_i \in [0,1]$ ở vòng $i$ (predictable bet).
    \item \textbf{Kết quả:} $R_i \in \{+1, -1\}$ (Đúng được $+1$, sai bị $-1$).
\end{itemize}

\vspace{0.3cm}
Quy tắc cập nhật tài sản sau mỗi vòng:
\[
L_i = L_{i-1}(1 + \lambda_i R_i)
\]

Sau $t$ vòng quan sát, tổng tài sản (\textit{wealth process}) là:
\[
L_t = \prod_{i=1}^{t} (1 + \lambda_i R_i)
\]

\end{frame}

%====================================================
% SAVI - Frame 13: Martingale dưới H0
%====================================================
\begin{frame}{Martingale dưới giả thuyết không}

Nếu giả thuyết không $H_0$ là đúng (Leila chỉ đoán mò):
\[
\mathbb{E}_{H_0}[R_i \mid \mathcal{F}_{i-1}] = 0
\]
Trong đó $\mathcal{F}_{i-1}$ là toàn bộ thông tin có sẵn trước vòng $i$.

Từ đó, kỳ vọng tài sản ở vòng tiếp theo là:
\[
\mathbb{E}_{H_0}[L_i \mid \mathcal{F}_{i-1}] = L_{i-1}
\]

\begin{block}{Tính chất Martingale}
Quá trình $(L_t)_{t\ge 0}$ là một \textbf{nonnegative martingale}. Dưới giả thuyết không, trò chơi là hoàn toàn công bằng.
\end{block}

\end{frame}

%====================================================
% SAVI - Frame 14: Định lý dừng và Bất đẳng thức
%====================================================
\begin{frame}{Optional Stopping \& Bất đẳng thức Ville}

Nhờ tính chất Martingale, ta có các kết quả lý thuyết mạnh mẽ:

\begin{itemize}
    \item \textbf{Định lý dừng tùy chọn (Optional Stopping Theorem):} Với mọi thời điểm dừng $\tau$,
    \[ \mathbb{E}_{H_0}[L_\tau] \le 1 \]
    \textit{(Khắc phục hoàn toàn vấn đề Optional Stopping của p-value!)}

    \vspace{0.2cm}
    \item \textbf{Bất đẳng thức Ville:} Phiên bản tuần tự của Bất đẳng thức Markov:
    \[ \Pr_{H_0} \left( \sup_{t\ge 0} L_t \ge \frac{1}{\alpha} \right) \le \alpha \]
    \textit{Ví dụ:} Với $\alpha = 0.02$, xác suất để tài sản $L_t \ge 50$ nhờ may mắn (nếu $H_0$ đúng) là không quá $2\%$.
\end{itemize}

\end{frame}

%====================================================
% SAVI - Frame 15: E-process và P-process
%====================================================
\begin{frame}{Suy luận tuần tự hợp lệ mọi thời điểm (SAVI)}

\begin{itemize}
    \item \textbf{e-process:} Quá trình tài sản $(L_t)_{t\ge 0}$ đóng vai trò đo lường bằng chứng. $L_t$ càng lớn, bằng chứng chống lại $H_0$ càng mạnh.
    \item \textbf{p-process:} Một anytime-valid p-value có thể được xây dựng từ e-process:
    \[ p_t = \inf_{s\le t} \frac{1}{L_s} \]
\end{itemize}

\begin{block}{Kiểm định tuần tự}
Ta bác bỏ giả thuyết không $H_0$ ngay khi tài sản vượt ngưỡng:
\[ \varphi_t = \mathbf{1} \left\{ L_t \ge \frac{1}{\alpha} \right\} \]
Nhờ BĐT Ville, xác suất sai lầm loại I được \textbf{kiểm soát chặt chẽ $\le \alpha$ tại mọi thời điểm dừng}.
\end{block}

\end{frame}

%====================================================
% SAVI - Frame 16: Tiêu chuẩn Kelly
%====================================================
\subsection{Cược Kelly và Tính Tối Ưu Logarit}

\begin{frame}{Tiêu chuẩn Kelly (Kelly Criterion)}

\begin{block}{Thiết lập bài toán}
Giả sử ta quan sát một chuỗi tung đồng xu độc lập (i.i.d.) với xác suất ngửa là $p > 1/2$ (giá trị $p$ đã biết).
\end{block}

\begin{itemize}
    \item \textbf{Vốn ban đầu:} \textbf{1} đô la.
    \item \textbf{Hành động:} Ở mỗi vòng, đặt cược một tỷ lệ $\lambda \in [0,1]$ tài sản vào mặt ngửa.
    \item \textbf{Luật chơi:} Ngửa thì nhân đôi tiền cược, sấp thì mất tiền cược.
\end{itemize}

\vspace{0.2cm}
Ký hiệu biến chỉ thị kết quả vòng thứ $i$:
\[
B_i =
\begin{cases}
1, & \text{nếu lần tung thứ } i \text{ là mặt ngửa},\\
-1, & \text{nếu lần tung thứ } i \text{ là mặt sấp}.
\end{cases}
\]

\end{frame}

%====================================================
% SAVI - Frame 17: Quá trình tài sản
%====================================================
\begin{frame}{Quá trình tài sản và Tốc độ tăng trưởng}

Tài sản sau $t$ vòng được biểu diễn bởi:
\[
W_t(\lambda) = \prod_{i=1}^{t} \left(1+\lambda B_i\right)
\]

Lấy logarit hai vế, ta chuyển phép nhân thành phép cộng:
\[
\log W_t(\lambda) = \sum_{i=1}^{t} \log\left(1+\lambda B_i\right)
\]

Theo luật số lớn (Law of Large Numbers), khi $t \to \infty$:
\[
\frac{1}{t}\log W_t(\lambda) \to \mathbb{E} \big[ \log(1+\lambda B) \big]
\]

\begin{alertblock}{Mục tiêu của Kelly}
Tìm tỷ lệ cược $\lambda$ để \textbf{tối đa hóa tốc độ tăng trưởng logarit dài hạn} của tài sản.
\end{alertblock}

\end{frame}

%====================================================
% SAVI - Frame 18: Tỷ lệ cược tối ưu
%====================================================
\begin{frame}{Nghiệm Kelly: Tỷ lệ cược tối ưu}

Ta cần giải bài toán tối ưu:
\[
\lambda^* = \arg\max_{\lambda} \mathbb{E} \big[ \log(1+\lambda B) \big]
\]

Khai triển hàm mục tiêu:
\[
\mathbb{E} \big[ \log(1+\lambda B) \big] = p\log(1+\lambda) + (1-p)\log(1-\lambda)
\]

Lấy đạo hàm theo $\lambda$ và cho bằng \textbf{0}:
\[
\frac{p}{1+\lambda} - \frac{1-p}{1-\lambda} = 0
\]

\begin{block}{Kết quả}
\[ \lambda^* = 2p - 1 = 2\left(p - \frac{1}{2}\right) \]
\textit{Ý nghĩa:} Tỷ lệ vốn đặt cược \textbf{tỷ lệ thuận} với mức độ lợi thế của người chơi so với một trò chơi công bằng ($p=1/2$). Nếu không có lợi thế ($p=1/2$), $\lambda^* = 0$ (không nên cược).
\end{block}

\end{frame}

%====================================================
% SAVI - Frame 19: Lý thuyết thông tin
%====================================================
\begin{frame}{Mối liên hệ với Lý thuyết Thông tin}

Thay $\lambda^*$ vào biểu thức tài sản, ta thu được tốc độ tăng trưởng tối ưu:
\[
W_t(\lambda^*) = \exp \Big( t\,H(p\|1/2) + o(t) \Big)
\]

Trong đó $H(p\|1/2)$ là \textbf{phân kỳ Kullback--Leibler} (Relative Entropy) giữa Bernoulli($p$) và Bernoulli($1/2$):
\[
H(p\|1/2) = p\log(2p) + (1-p)\log\big(2(1-p)\big)
\]

Tương đương:
\[
\lim_{t\to\infty} \frac{\mathbb{E}[\log W_t(\lambda^*)]}{t} = H(p\|1/2)
\]

\begin{alertblock}{Kết luận}
Tốc độ tăng trưởng tối ưu của tài sản chính là \textbf{lượng thông tin} mà phân phối thật chứa đựng so với giả thuyết đồng xu công bằng.
\end{alertblock}

\end{frame}

%====================================================
% SAVI - Frame 20: Vai trò trong SAVI
%====================================================
\begin{frame}{Tầm quan trọng của Cược Kelly}

Bên cạnh việc tối ưu tốc độ tăng trưởng dài hạn, tiêu chuẩn Kelly còn mang lại nhiều đặc tính tối ưu tiệm cận khác:

\begin{itemize}
    \item \textbf{Tối thiểu hóa thời gian kỳ vọng} để đạt tới một ngưỡng tài sản cho trước.
    \item \textbf{Tối đa hóa tài sản kỳ vọng} tại một thời điểm xác định.
\end{itemize}

\vspace{0.3cm}
\begin{block}{Viên gạch nền tảng của Game-Theoretic Statistics}
Chính sự kết nối bền chặt giữa \textit{Lý thuyết cá cược} và \textit{Lý thuyết thông tin} đã biến Tiêu chuẩn Kelly thành nền tảng để xây dựng các \textbf{test martingale} và \textbf{e-process} trong phương pháp suy luận tuần tự \textbf{SAVI}.
\end{block}

\end{frame}
